\documentclass[11pt,a4paper]{article}

\usepackage[T1]{fontenc}
\usepackage[utf8]{inputenc}
\usepackage{lmodern}
\usepackage[english]{babel}
\usepackage{microtype}
\usepackage{geometry}
\geometry{margin=2.4cm}

\usepackage{amsmath,amssymb,amsthm,mathtools,bm}
\usepackage{enumitem}
\usepackage{graphicx}
\usepackage{algorithm}
\usepackage{algpseudocode}
\usepackage{booktabs}
\usepackage{xcolor}
\usepackage[numbers,sort&compress]{natbib}
\usepackage[colorlinks,citecolor=blue!50!black,linkcolor=blue!50!black]{hyperref}
\usepackage[nameinlink,noabbrev]{cleveref}

\newcommand{\eqjust}[1]{\qquad\eqref{#1}}

\theoremstyle{plain}
\newtheorem{theorem}{Theorem}
\newtheorem{lemma}{Lemma}
\newtheorem{proposition}{Proposition}
\theoremstyle{definition}
\newtheorem{example}{Example}
\theoremstyle{remark}
\newtheorem{remark}{Remark}

\newcommand*\R{\mathbb{R}}
\newcommand*\N{\mathbb{N}}
\newcommand*\x{\mathbf{x}}
\newcommand*\U{\mathbf{u}}
\newcommand*\y{\mathbf{y}}
\newcommand*\z{\mathbf{z}}
\newcommand*\V{\mathbf{v}}
\newcommand*\bv{\mathbf{b}}
\newcommand*\A{\mathbf{A}}
\newcommand*\prox{\operatorname{prox}}
\DeclareMathOperator*{\argmin}{arg\,min}
\newcommand*{\norm}[1]{\|#1\|}

\algnewcommand{\Input}[1]{\textbf{Input:} #1}
\algnewcommand{\Output}[1]{\textbf{Output:} #1}

\title{The group-sparse proximal mapping\\with lower and upper cardinality bounds}
\author{Gabriel Belém Barbosa\thanks{IMECC, University of Campinas (UNICAMP).
\texttt{gggt4@ime.unicamp.br}}
\and Paulo J.\ S.\ Silva\thanks{IMECC, University of Campinas (UNICAMP).}}
\date{September 2026}

\begin{document}
\maketitle

\begin{abstract}
We study composite minimization of a smooth objective plus a group-\(\ell_0\)
regularizer, subject to both an upper and a \emph{lower} bound on the number of
active groups. The lower bound models diversification or coverage constraints
that the classical cardinality-constrained formulation (\(l=0\)) cannot
express. The main result is a complete characterization of the proximal mapping
of the resulting discontinuous regularizer, extending
\cite[Theorem~3.2]{beck2018}. In particular, we identify the set of
activatable groups, prove that \(\prox_h\) is empty precisely when fewer than \(l\) groups can
be activated, and give an explicit description of every proximal point together
with an algorithm that computes one. The argument follows the comparison
strategy of \cite{beck2018} for \(\omega_{[l]}>2\lambda\), and treats the new case \(\omega_{[l]}\le 2\lambda\) by a separate sandwich of
supports and an explicit reconstruction of the witnessing index sets.
\end{abstract}

\noindent\textbf{Keywords:}
group sparsity; \(\ell_0\) regularization; proximal mapping; cardinality
constraints; composite optimization.

\section{Introduction}
\label{sec:intro}

Group sparsity comprises optimization problems whose decision variable is
partitioned into disjoint blocks (groups), and in which one wishes to limit or
penalize the number of nonzero groups. The structure arises whenever
explanatory variables are highly correlated and must be treated jointly, or
when the model itself imposes a predetermined grouping
\cite{yuan2006,simon2013,beck2018}. Standard applications include
portfolio construction,
imaging and signal processing, total-variation denoising, compressed sensing
\cite{donoho2006}, and low-rank approximation; see \cite{beck2018} and the
references therein.

The generic proximal mapping of \cite{beck2018} covers a penalty on the number
of active groups together with an \emph{upper} bound \(s\) on that number. We enlarge the class slightly by introducing a \emph{lower}
bound \(l\) as well. The extra constraint controls the maximal sparsity of a
solution and is useful whenever a minimum level of diversification is
required: one should not concentrate a portfolio on a handful of sectors
\cite{brodie2009,beck2018}. Even when the floor is only a target rather
than a hard modelling constraint, specifying \(l\) is more convenient than
tuning \(\lambda\) until the recovered support meets a desired minimum size:
the constraint enforces the cardinality directly, and the grid of candidate
penalties can be cut accordingly. The aim of this paper is a closed form for the proximal calculus
of the resulting regularizer and that can be used as a black box inside
forward--backward methods \cite{moreau1962,beckbook}.

\paragraph{Contributions.}
We characterize \(\prox_h\) for
\[
h(\x)=\lambda\|g(\x)\|_0+\delta_{B\cap C_{l,s}}(\x),
\]
where \(C_{l,s}\) is the set of vectors with between \(l\) and \(s\) active
groups and \(B\) is a product of closed group-wise constraints (each group may
be set to the origin). The new analytical points are:
\begin{itemize}[leftmargin=1.4em,itemsep=0.25em]
\item the \emph{activatable} index set \(I_{\mathcal{A}}(\x)\), on which all
order statistics of the Beck--Hallak gains \(\omega(\x)\) are computed;
\item emptiness of \(\prox_h\) if and only if \(l\ge 1\) and
\(|I_{\mathcal{A}}(\x)|<l\);
\item a two-case description of the active support of a proximal solution, according to the position
of the \(l\)-th order statistic \(\omega(\x)_{[l]}\) relative to the threshold
\(2\lambda\);
\item a closed-form family \(\mathcal{T}(\x)\) of admissible supports and an
algorithm that returns one proximal point (or reports emptiness).
\end{itemize}
When \(\omega(\x)_{[l]}>2\lambda\) the argument is conceptually the same as
\cite[Theorem~3.2]{beck2018}. When \(\omega(\x)_{[l]}\le 2\lambda\) one must
force additional zero or negative gain groups in order to meet the lower bound \(l\); that
case has no counterpart for \(l=0\).

\section{Problem formulation}
\label{sec:problem}

Fix a partition \(\{G_i\}_{i=1}^m\) of \(\{1,2,\ldots,n\}\) into \(m\) groups
of sizes \(n_1,\ldots,n_m\). After a harmless reindexing,
\[
\begin{aligned}
G_1 & \coloneqq\left\{1,2, \ldots, n_1\right\}, G_2\coloneqq\left\{n_1+1, n_1+2, \ldots, n_1+n_2\right\}, \ldots, \\
G_m & \coloneqq\left\{n_{m-1}+1, n_{m-1}+2, \ldots, n\right\} .
\end{aligned}
\]
For \(\x\in\R^n\) and \(T\subset \{1,\ldots,n\}\), write \(\x_T\) for the
subvector of \(\x\) indexed by \(T\). The map \(g: \R^n \to\{0,1\}^m\) that
marks nonempty groups is
$$
g(\x)_i\coloneqq \begin{cases}1, & \x_{G_i} \neq \mathbf{0}, \\
0, & \text{otherwise }\end{cases}.
$$
Thus \(g(\x)_i=1\) means that group \(G_i\) is \emph{active} for \(\x\);
otherwise it is \emph{inactive}. For \(l,s\in \N_0\) with \(l\leq s\leq m\),
$$
C_{l,s}\coloneqq\left\{\x \in \R^n:l\leq\|g(\x)\|_0 \leq s\right\}
$$
is the set of vectors with at least \(l\) and at most \(s\) active groups.
Ordinary (ungrouped) sparsity is the special case in which every group is a
singleton, so that \(C_{l,s}\) consists of the \(s\)-sparse vectors with at
least \(l\) nonzero entries.

The problem under study is
\begin{equation}
\tag{P0}
\label{P0}
\begin{aligned}
&\min&& f(\x)+\lambda\|g(\x)\|_0 \\
&\operatorname{s.t.}&& \x \in C_{l,s} \cap B,
\end{aligned}
\end{equation}
where \(\lambda \geq 0\) is a penalty parameter, \(f: \R^n \to \R\) is
continuously differentiable and bounded from below, and
\begin{equation}
\label{B}
B \coloneqq \prod_{i=1}^m\left(D_i \cup\{\mathbf{0}\}\right)=\left(D_1 \cup\{\mathbf{0}\}\right) \times\left(D_2 \cup\{\mathbf{0}\}\right) \times \cdots \times\left(D_m \cup\{\mathbf{0}\}\right),
\end{equation}
with each \(D_i \subset \R^{n_i}\) nonempty and closed. This problem is equivalent, in
composite unconstrained form,
\begin{equation}
\label{P}
\min _{\x \in \R^n} f(\x)+h(\x),
\end{equation}
with
\begin{equation}
\label{h}
h(\x) \coloneqq\lambda\|g(\x)\|_0+\delta_{B \cap C_{l,s}}(\x).
\end{equation}
Problem~\eqref{P} has the familiar shape of proximal composite optimization.
The proximal mapping of \(h\)
\begin{equation}
\label{prox}
\begin{aligned}
\prox_{h}(\x)&\coloneqq\argmin_{\U\in\R^n}\left\{h(\U)+ \frac{1}{2}\norm{\U-\x}^2 \right\}
\end{aligned}
\end{equation}
was characterized in \cite{beck2018} for \(l=0\). \Cref{thm:proxh} below
extends that characterization to \(l\leq s\).

\section{Motivating examples}
\label{sec:examples}

\begin{example}[sector-constrained portfolio]
\label{ex:carteira}
Following \cite{beck2018}, consider mean--variance allocation with a bound on
the number of sectors that may receive capital,
\begin{equation}
\begin{aligned}
&\min&& \gamma \cdot \x^T \mathbf{C} \x-\boldsymbol{\mu}^T \x \\
&\operatorname{s.t.}&& \x \in C_{l,s} \cap \Delta_n
\end{aligned}.
\end{equation}
Here \(\boldsymbol{\mu} \in \R^n\) is a vector of expected returns,
\(\mathbf{C} \in \R^{n \times n}\) is a covariance matrix, \(\gamma>0\) trades
off risk against return, and
$$\Delta_n=\left\{\x \in \R^n\mid \x \geq \mathbf{0},\ \sum_{i=1}^n \x_i=1\right\}$$
is the unit simplex. Decision variables are already grouped by sector. An
upper bound \(s\) limits concentration; a positive lower bound \(l\)
enforces a minimal diversification
across sectors, in the same spirit as sparse Markowitz models
\cite{brodie2009}.
\end{example}

\begin{example}[cardinality-constrained sparse regression]
\label{ex:imbuia}
An instance of \eqref{P} with singleton groups is
\begin{equation}
\label{regress2}
\min_{\x\in\R^n}\frac{1}{2}\|\A\x-\bv\|^2+\lambda\|\x\|_0+\delta_{C_{l,s}}(\x),
\end{equation}
with \(\A\in \R^{m\times n}\) and \(\bv\in \R^m\).
The formulation was motivated by infrared spectroscopy on the IMBUIA
beamline (SIRIUS/LNLS), where a sparse spectrogram is recovered from
measurements.
If a prior estimate of the number of spectral lines is available --- from
previous experiments or from the literature --- it can be encoded as the
lower bound \(l\), so that the \(\ell_0\) penalty is not free to discard
known structure.
\end{example}

\section{Notation}
\label{sec:notation}

We first present a few auxiliary objects. Given \(\V \in \R^m\) and
\(i=1,2,\ldots,m\), \(\V_{[i]}\) denotes the \(i\)-th largest entry of
\(\V\), so that \(\V_{[1]}\geq \V_{[2]}\geq\hdots\geq\V_{[m]}\). The family
\begin{equation}\label{defS}
    \mathcal{S}_i(\V)\coloneqq\{S \cup T\mid S\coloneqq\{j\mid\V_j>\V_{[i]}\},\ T\subset\{j\mid\V_j=\V_{[i]}\},\ | S \cup T | = i\}\quad \forall i=1,2,\ldots,m
 \end{equation}
collects every index set of cardinality \(i\) associated with the largest
entries of \(\V\). The set \(\mathcal{S}_i(\V)\) need not be a singleton,
since \(\V\) may have ties. For instance, if \(\V = (0\ 0\ 1)^T\), then
\(\mathcal{S}_2(\V) = \{\{1,3\},\{2,3\}\}\), where
\(S=\{j\mid\V_j>\V_{[2]}\}=\{3\}\) and the admissible subsets of
\(\{j\mid\V_j=\V_{[2]}\}\) are \(T=\{1\}\) and \(T=\{2\}\).

Given \(\x \in \R^n\), the orthogonal projection onto a set \(C\subset\R^n\)
is
$$
P_C(\x)\coloneqq \prox_{\delta_C}(\x)=\argmin_{\y\in C}\|\y-\x\|_2,
$$
and the distance to \(D_i\) is
$$
d_{D_i}\left(\x_{G_i}\right)\coloneqq\min _{\z \in D_i}\left\|\x_{G_i}-\z\right\|_2\quad \forall i=1,2, \ldots, m .
$$
The latter is well defined because each \(D_i\) is closed. The gain map
\(\omega: \R^n \to \R^m\),
\begin{equation}
\label{omega}
\omega(\x)_i\coloneqq\left\|\x_{G_i}\right\|_2^2-d_{D_i}^2\left(\x_{G_i}\right)\quad \forall i=1,2, \ldots, m,
\end{equation}
governs which groups are activated by \(\prox_h\).

The support of a vector (its set of active-group indices) is
$$I_1(\x)\coloneqq\left\{i \in\{1,2, \ldots, m\}\mid g(\x)_i=1\right\}.$$
The origin is always feasible for every group, through
\(B=\prod_i(D_i\cup\{\mathbf{0}\})\). Activating group \(i\), however,
requires a point of \(D_i\setminus\{\mathbf{0}\}\). This motivates the set of
\emph{activatable} groups
\begin{equation}
\label{Iativ}
I_{\mathcal{A}}(\x)\coloneqq\bigl\{i\in\{1,2,\ldots,m\}\mid P_{D_i}(\x_{G_i})\setminus\{\mathbf{0}\}\neq\emptyset\bigr\}.
\end{equation}
If \(\mathbf{0}\) is an accumulation point of \(D_i\) and the unique nearest
point is the origin---for example \(D_1=[0,1]\) and \(\x=\mathbf{0}\)---then
\(i\notin I_{\mathcal{A}}(\x)\): the minimum over \(D_i\setminus\{\mathbf{0}\}\)
is not attained, so that group cannot contribute a block of a proximal point.

\begin{lemma}\label{lem:ativavel}
Let \(\x\in\R^n\) and \(i\in\{1,2,\ldots,m\}\). If \(\omega(\x)_i\neq 0\), then
\(i\in I_{\mathcal{A}}(\x)\).
\end{lemma}
\begin{proof}
Since \(D_i\) is closed, \(P_{D_i}(\x_{G_i})\) is nonempty and \(\omega(\x)_i\) well defined. If \(\omega(\x)_i>0\), then \(d_{D_i}(\x_{G_i})<\|\x_{G_i}\|_2\), hence
\(\mathbf{0}\notin P_{D_i}(\x_{G_i})\), therefore
\(P_{D_i}(\x_{G_i})\setminus\{\mathbf{0}\}\neq\emptyset\). If
\(\omega(\x)_i<0\), then \(d_{D_i}(\x_{G_i})>\|\x_{G_i}\|_2\), which forces
\(\mathbf{0}\notin D_i\). Thus \(P_{D_i}(\x_{G_i})\subset D_i\) does not
contain the origin.
\end{proof}
A point of \(D_i\) strictly closer than the origin would yield
\(\omega(\x)_i>0\). Consequently, if \(\omega(\x)_i=0\), the origin already
realizes the distance to \(D_i\cup\{\mathbf{0}\}\); one has
\(i\in I_{\mathcal{A}}(\x)\) only when \(P_{D_i}(\x_{G_i})\) contains another
point \emph{at the same} distance.

When applied to \(\omega(\x)\), the order statistics and the sets
are taken over \(I_{\mathcal{A}}(\x)\): \(\omega(\x)_{[k]}\) is
the \(k\)-th largest value in
\(\{\omega(\x)_i\mid i\in I_{\mathcal{A}}(\x)\}\), and
\(\mathcal{S}_k(\omega(\x))\) is \eqref{defS} for that subvector. Likewise,
$$I_{a}^{\stackrel{?}{=}}(\x)\coloneqq\left\{i\in I_{\mathcal{A}}(\x)\mid \omega(\x)_i\stackrel{?}{=}a\right\},$$
where \(\stackrel{?}{=}\) stands for an arbitrary relation, and \(a\in\R\) may
also be written as \([i]\) to mean \(\omega(\x)_{[i]}\). We write
\(s'\coloneqq\min\{s,|I_{\mathcal{A}}(\x)|\}\).

For an index set \(T\subset\{1,2,\ldots,n\}\), the matrix \(\mathbf{U}_T\)
consists of the columns of \(\mathbf{I}_n\) indexed by \(T\), and we write
$$
\mathcal{A}(T)\coloneqq\bigcup_{i\in T} G_i
$$
and
$$
\mathcal{B}(T)\coloneqq \prod_{i\in T}D_i.
$$

The following identity is \cite[Lemma~3.1]{beck2018}.

\begin{lemma}\label{lem:toomega}
Let \(\x \in \R^n\) and \(T, S \subset\{1,2, \ldots, m\}\). Then, for all
\(\z, \y \in \R^n\) satisfying
$$
\z_{G_i} \in\left\{\begin{array} { l l }
{ P _ { D _ { i } } ( \x _ { G _ { i } } ) , } & { i \in T  } \\
{ \{ \mathbf { 0 } \} , } & { i \notin T  }
\end{array}, \quad \y _ { G _ { i } } \in \left\{\begin{array}{ll}
P_{D_i}\left(\x_{G_i}\right), & i \in S \\
\{\mathbf{0}\}, & i \notin S
\end{array}\right.\right.,
$$
one has
$$
\|\z-\x\|_2^2-\|\y-\x\|_2^2=\sum_{i \in S} \omega(\x)_i-\sum_{i \in T} \omega(\x)_i .
$$
\end{lemma}

\section{When the proximal mapping is empty}
\label{sec:empty}

\begin{proposition}\label{prop:empty}
One has \(\prox_h(\x)=\emptyset\) if and only if \(l\geq 1\) and
\(|I_{\mathcal{A}}(\x)|<l\).
\end{proposition}
\begin{proof}
If \(\U\in\prox_h(\x)\), then for
\(i\in I_1(\U)\) one has \(\U_{G_i}\neq\mathbf{0}\) and \(\U_{G_i}\in P_{D_i}(\x_{G_i})\), otherwise a better solution of the problem defining \(\prox_h(\x)\) could be
obtained with a smaller quadratic term and the same active groups. A group with
\(P_{D_i}(\x_{G_i})=\{\mathbf{0}\}\) cannot be activated, thus \(I_1(\U)\subset I_{\mathcal{A}}(\x)\). Since \(|I_1(\U)|\geq l\), it follows that
\(|I_{\mathcal{A}}(\x)|\geq l\).

Conversely, if \(|I_{\mathcal{A}}(\x)|\geq l\) or \(l=0\),
take \(T\subset I_{\mathcal{A}}(\x)\) maximizing
\(\sum_{i\in T}\bigl(\omega(\x)_i-2\lambda\bigr)\) subject to
\(l\leq |T|\leq s'\) (so \(T=\emptyset\) is admissible when \(l=0\)), and set
\(\U_{G_i}\in P_{D_i}(\x_{G_i})\setminus\{\mathbf{0}\}\) for \(i\in T\) and
\(\U_{G_i}=\mathbf{0}\) otherwise. Then \(\U\in B\cap C_{l,s}\) and, by
\Cref{lem:toomega},
\[
\|\U-\x\|_2^2+2\lambda\|g(\U)\|_0=\|\x\|_2^2-\sum_{i\in T}\bigl(\omega(\x)_i-2\lambda\bigr).
\]
Any \(\y\in B\cap C_{l,s}\) with \(S=I_1(\y)\subset I_{\mathcal{A}}(\x)\)
satisfies
\[
\|\y-\x\|_2^2+2\lambda\|g(\y)\|_0\geq\|\x\|_2^2-\sum_{i\in S}\bigl(\omega(\x)_i-2\lambda\bigr),
\]
because \(\y_{G_i}\in D_i\setminus\{\mathbf{0}\}\) for \(i\in S\), and
equality in \Cref{lem:toomega} requires the projections. The set \(S\) also
obeys \(l\leq |S|\leq s'\), hence, by construction,
\[
\sum_{i\in T}\bigl(\omega(\x)_i-2\lambda\bigr)\geq\sum_{i\in S}\bigl(\omega(\x)_i-2\lambda\bigr).
\]
A support leaving \(I_{\mathcal{A}}(\x)\) does not attain the minimum.
Therefore \(\U\in\prox_h(\x)\).
\end{proof}

\section{Characterization of the proximal mapping}
\label{sec:prox}

We now characterize \(\prox_{h}\) for variable \(l\). The argument follows
\cite[Theorem~3.2]{beck2018} closely: the same comparison of gains is used,
and the first half of the proof, corresponding to \(\omega(\x)_{[l]}> 2\lambda\),
is conceptually identical to the \(l=0\) case, up to a few details.

\begin{theorem}\label{thm:proxh}
Let \(\x \in \R^n\).
If \(|I_{\mathcal{A}}(\x)|<l\), then \(\prox_h(\x)=\emptyset\). If
\(|I_{\mathcal{A}}(\x)|\geq l\), then \(\U \in \prox_h(\x)\) if and only if
the following conditions hold:
\begin{enumerate}[label=(\alph*)]
    \item \(\U_{G_i} \in P_{D_i}\left(\x_{G_i}\right)\ \forall i \in I_1(\U)\);
    \item If \(\omega(\x)_{[l]}> 2\lambda\) or \(l=0\), \(\exists T \in \mathcal{S}_{s'}(\omega(\x))\) such that
\begin{equation}
\label{T1}
    T \cap I_{2\lambda}^>(\x) \subset I_1(\U) \subset T \cap I_{2\lambda}^\geq(\x),
\end{equation}
and if \(\omega(\x)_{[l]}\leq 2\lambda\) and \(l\geq 1\), \(\exists T \in \mathcal{S}_{l-1}(\omega(\x))\), \(i\in I^=_{[l]}(\x)\setminus T\) and
$$
U\in
\begin{cases}
    \left\{W \subset I_{2 \lambda}^{=}(x) \backslash(T \cup\{i\})\mid |W| =\min \bigl\{\bigl|I_{2 \lambda}^=(\x)\setminus(T\cup\{i\})\bigr|, s'-l\bigr\}\right\},&\omega(\x)_{[l]}=2\lambda\\
    \{\emptyset\},&\omega(\x)_{[l]}<2\lambda
\end{cases}
$$
such that
\begin{equation}
\label{T2}
T\cup \{i\}\subset I_1(\U)\subset T\cup\{i\}\cup U.
\end{equation}
\end{enumerate}
\end{theorem}
\begin{proof}
The case \(|I_{\mathcal{A}}(\x)|<l\) is \Cref{prop:empty}. From now on,
\(|I_{\mathcal{A}}(\x)|\geq l\).
Suppose that \(\U\) satisfies (a) and (b), and let
\(\y \in \prox_{h}(\x)\).
We will show that \(\U \in \prox_h(\x)\) and that \(\y\)
satisfies (a) and \eqref{T1} or \eqref{T2} (as appropriate) for some
\(\tilde{T}\) and, if needed, \(\tilde{U}\), defined as in (b). Note that, by
(a), \(\U\in B\), and by (b), \(\U\in C_{l,s}\) (this will be checked in more
detail for each case). Therefore, since \(\y \in \prox_h(\x)\)
(hence also \(\y\in B\cap C_{l,s}\)),
\begin{equation}
\label{proxy}
    2 \lambda\left(\|g(\U)\|_0-\|g(\y)\|_0\right)+\|\U-\x\|_2^2-\|\y-\x\|_2^2 \geq 0 .
\end{equation}
Plainly, \(\y_{G_i} \in P_{D_i}\left(\x_{G_i}\right)\) \(\forall i \in I_1(\y)\)
(otherwise a better solution of the problem defining
\(\prox_h(\x)\) could be obtained with the same active groups)
and, by hypothesis, \(\U_i \in P_{D_i}\left(\x_{G_i}\right)\ \forall i \in I_1(\U)\).
Hence, by \Cref{lem:toomega}, \eqref{proxy} is the same as
$$
\sum_{i \in I_1(\y)}\left(\omega(\x)_i-2 \lambda\right)-\sum_{i \in I_1(\U)}\left(\omega(\x)_i-2 \lambda\right) \geq 0,
$$
which in turn is equivalent to
\begin{equation}
\label{main1}
\begin{aligned}
&\sum_{i \in I_1(\y) \cap I_{2\lambda}^>(\x)}\left(\omega(\x)_i-2 \lambda\right)+\sum_{i \in I_1(\y) \cap I^<_{2\lambda}(\x)}\left(\omega(\x)_i-2 \lambda\right)\\
&\qquad-\sum_{i \in I_1(\U)}\left(\omega(\x)_i-2 \lambda\right) \geq 0 .
\end{aligned}
\end{equation}

We first treat the first case of (b). The case \(l=0\) is
\cite[Theorem~3.2]{beck2018}. If \(l\geq 1\), \eqref{T1} indeed yields
\(l\leq |I_1(\U)|\leq s'\), since \(|T \cap I_{2\lambda}^>(\x)|\geq l\)
because \(\omega(\x)_{[l]}> 2\lambda\), and \(|T|=s'\) by \eqref{defS}.

Any indices of \(I_1(\U)\) that lie in \(I_{2\lambda}^=(\x)\) do not change the sum, hence by \eqref{T1}
$$
\sum_{i \in I_1(\U)}\left(\omega(\x)_i-2 \lambda\right)=\sum_{i \in T \cap I_{2\lambda}^\geq(\x)}\left(\omega(\x)_i-2 \lambda\right)=\sum_{i \in T \cap I_{2\lambda}^>(\x)}\left(\omega(\x)_i-2 \lambda\right) .
$$
Since \(T\) contains indices of a set of the \(s'\) largest entries of
\(\omega(\x)\) and \(\y\) has at most \(s\) active groups (because
\(\y \in C_{l,s}\)), one obtains
$$
\sum_{i \in T \cap I_{2\lambda}^\geq(\x)}\left(\omega(\x)_i-2 \lambda\right) \geq \sum_{i \in I_1(\y) \cap I_{2\lambda}^>(\x)}\left(\omega(\x)_i-2 \lambda\right) .
$$
Combining the last two relations yields
\begin{equation}
\label{ubest}
\sum_{i \in I_1(\U)}\left(\omega(\x)_i-2 \lambda\right) \geq \sum_{i \in I_1(\y) \cap I_{2\lambda}^>(\x)}\left(\omega(\x)_i-2 \lambda\right) .
\end{equation}

Using the elementary inequality
\begin{equation}
\label{negat}
    \sum_{i \in I_1(\y) \cap I^<_{2\lambda}(\x)}\left(\omega(\x)_i-2 \lambda\right) \leq 0
\end{equation}
together with \eqref{ubest} gives
\begin{equation}
    \label{desig}
\begin{aligned}
0 & \geq \sum_{i \in I_1(\y) \cap I_{2\lambda}^>(\x)}\left(\omega(\x)_i-2 \lambda\right)-\sum_{i \in I_1(\U)}\left(\omega(\x)_i-2 \lambda\right) \\
& \geq \sum_{i \in I_1(\y) \cap I_{2\lambda}^>(\x)}\left(\omega(\x)_i-2 \lambda\right)+\sum_{i \in I_1(\y) \cap I^<_{2\lambda}(\x)}\left(\omega(\x)_i-2 \lambda\right) \\
& \qquad-\sum_{i \in I_1(\U)}\left(\omega(\x)_i-2 \lambda\right) \\
& \geq 0. \eqjust{main1}
\end{aligned}
\end{equation}
Thus the chain \eqref{desig} holds as a chain of equalities, and therefore
\eqref{proxy} holds as an equality as well, which implies
\(\U \in \prox_h(\x)\) by the definition \eqref{prox}. The fact
that \eqref{desig} is a chain of equalities also implies
$$
\sum_{i \in I_1(\y) \cap I_{2\lambda}^>(\x)}\left(\omega(\x)_i-2 \lambda\right)=\sum_{i \in I_1(\U)}\left(\omega(\x)_i-2 \lambda\right)
$$
and
$$
\sum_{i \in I_1(\y) \cap I^<_{2\lambda}(\x)}\left(\omega(\x)_i-2 \lambda\right)=0,
$$
which is equivalent to \(I_1(\y) \cap I^<_{2\lambda}(\x)=\emptyset\) by
negativity of the summands. The last two relations in turn imply
$$
\sum_{i \in I_1(\y)}\left(\omega(\x)_i-2 \lambda\right)=\sum_{i \in I_1(\y) \cap I_{2\lambda}^>(\x)}\left(\omega(\x)_i-2 \lambda\right)=\sum_{i \in I_1(\U)}\left(\omega(\x)_i-2 \lambda\right) .
$$

It remains to exhibit \(\tilde{T}\in\mathcal{S}_{s'}(\omega(\x))\) such that
\eqref{T1} holds with \(\U\) replaced by \(\y\).
One already has \(I_1(\y)\cap I_{2\lambda}^<(\x)=\emptyset\),
\(l\leq |I_1(\y)|\leq s'\)
and by the above equation
\[
\sum_{i\in I_1(\y)\cap I_{2\lambda}^>(\x)}\bigl(\omega(\x)_i-2\lambda\bigr)=\sum_{i\in T\cap I_{2\lambda}^>(\x)}\bigl(\omega(\x)_i-2\lambda\bigr).
\]
The terms in \(I_{2\lambda}^>(\x)\) are strictly positive and those in
\(I_{2\lambda}^=(\x)\) vanish.

If \(\omega(\x)_{[s']}>2\lambda\), then \(T\cap I_{2\lambda}^>(\x)=T\) has
cardinality \(s'\) and realizes the maximal sum among subsets of \(s'\)
indices. A proper subset of \(I_{2\lambda}^>(\x)\), or the swap of an index
in \(I_{2\lambda}^>(\x)\) for one in \(I_{2\lambda}^=(\x)\), strictly
decreases the sum. Hence \(I_1(\y)\cap I_{2\lambda}^>(\x)\in\mathcal{S}_{s'}(\omega(\x))\)
and \(|I_1(\y)\cap I_{2\lambda}^>(\x)|=s'\), which forces
\(I_1(\y)\subset I_{2\lambda}^>(\x)\). Take \(\tilde{T}\coloneqq I_1(\y)\). Then
\(\tilde{T}\in\mathcal{S}_{s'}(\omega(\x))\) and
\(\tilde{T}\cap I_{2\lambda}^>(\x)=I_1(\y)=\tilde{T}\cap I_{2\lambda}^\geq(\x)\).

If \(\omega(\x)_{[s']}=2\lambda\), every set in \(\mathcal{S}_{s'}(\omega(\x))\)
contains \(I_{2\lambda}^>(\x)\), hence
\(T\cap I_{2\lambda}^>(\x)=I_{2\lambda}^>(\x)\). Equality of the positive
sums forces \(I_1(\y)\cap I_{2\lambda}^>(\x)=I_{2\lambda}^>(\x)\). Together with
\(I_1(\y)\cap I_{2\lambda}^<(\x)=\emptyset\), this yields
\(I_{2\lambda}^>(\x)\subset I_1(\y)\subset I_{2\lambda}^\geq(\x)\). Since
\(|I_1(\y)|\leq s'\) and \(|I_{2\lambda}^\geq(\x)|\geq s'\), there exists
\(W\subset I_{2\lambda}^=(\x)\) with \(I_1(\y)\setminus I_{2\lambda}^>(\x)\subset W\)
and \(|W|=s'-|I_{2\lambda}^>(\x)|\). Take
\(\tilde{T}\coloneqq I_{2\lambda}^>(\x)\cup W\). Then
\(\tilde{T}\in\mathcal{S}_{s'}(\omega(\x))\) and
\(\tilde{T}\cap I_{2\lambda}^>(\x)\subset I_1(\y)\subset \tilde{T}\cap I_{2\lambda}^\geq(\x)\).

If \(\omega(\x)_{[s']}<2\lambda\), again
\(T\cap I_{2\lambda}^>(\x)=I_{2\lambda}^>(\x)\) and equality of the positive
sums gives \(I_{2\lambda}^>(\x)\subset I_1(\y)\subset I_{2\lambda}^\geq(\x)\). Any
\(\tilde{T}\in\mathcal{S}_{s'}(\omega(\x))\) satisfies
\(\tilde{T}\cap I_{2\lambda}^>(\x)=I_{2\lambda}^>(\x)\) and
\(\tilde{T}\cap I_{2\lambda}^\geq(\x)=I_{2\lambda}^\geq(\x)\), so \eqref{T1}
holds for \(\y\).

In all subcases, \(|I_1(\y)|\geq l\): if \(\omega(\x)_{[s']}>2\lambda\) one has
\(|I_1(\y)|=s'\geq l\); in the others, \(I_{2\lambda}^>(\x)\subset I_1(\y)\) and
\(|I_{2\lambda}^>(\x)|\geq l\) because \(\omega(\x)_{[l]}>2\lambda\) (or
\(l=0\)).
Finally, as already observed, \(\y_{G_i} \in P_{D_i}\left(\x_{G_i}\right)\)
\(\forall i \in I_1(\y)\), so (a) also holds for \(\y\).

We now treat the second case of (b), in which \(\omega(\x)_{[l]}\leq 2\lambda\).
Indeed, \eqref{T2} gives \(l\leq |I_1(\U)|\leq s'\), since
\(|T\cup \{i\}|=l\) by \eqref{defS}, and \(|T\cup\{i\}\cup U|\leq s'\) because
\(|U|\leq s'-l\) by construction of \(U\).

If \(\omega(\x)_{[l]}=2\lambda\), one has \(I_{[l]}^>(\x)=I_{2\lambda}^>(\x)\)
and \(|I_{2\lambda}^>(\x)|\leq l-1\). Since \(T\in\mathcal{S}_{l-1}(\omega(\x))\),
every index of \(I_{2\lambda}^>(\x)\) is among the \(l-1\) largest, hence
\(I_{2\lambda}^>(\x)\subset T\). By \eqref{T2},
\(I_1(\U)\setminus I_{2\lambda}^>(\x)\subset I_{2\lambda}^=(\x)\), and therefore
\(I_{2\lambda}^>(\x)\subset I_1(\U)\subset I_{2\lambda}^\geq(\x)\). The terms
in \(I_{2\lambda}^=(\x)\) vanish, whence
\[
\sum_{i\in I_1(\U)}\bigl(\omega(\x)_i-2\lambda\bigr)=\sum_{i\in I_{2\lambda}^>(\x)}\bigl(\omega(\x)_i-2\lambda\bigr)\geq\sum_{i\in I_1(\y)\cap I_{2\lambda}^>(\x)}\bigl(\omega(\x)_i-2\lambda\bigr).
\]
Together with \eqref{negat} and \eqref{main1}, the chain \eqref{desig} holds
again and is satisfied as equality. Thus \(\U\in\prox_h(\x)\),
\(I_1(\y)\cap I_{2\lambda}^<(\x)=\emptyset\), and equality of the positive
sums forces \(I_{2\lambda}^>(\x)\subset I_1(\y)\). One obtains
\(I_{2\lambda}^>(\x)\subset I_1(\y)\subset I_{2\lambda}^\geq(\x)\) and, since
\(\y\in C_{l,s}\), \(l\leq |I_1(\y)|\leq s'\).

It remains to exhibit \(\tilde{T}\in\mathcal{S}_{l-1}(\omega(\x))\),
\(\tilde{\imath}\in I^=_{[l]}(\x)\setminus\tilde{T}\) and \(\tilde{U}\) as in
\eqref{T2} such that \eqref{T2} holds with \(\U\) replaced by \(\y\). Since
\(|I_1(\y)|\geq l>|I_{2\lambda}^>(\x)|\), the set \(I_1(\y)\cap I_{2\lambda}^=(\x)\) is
nonempty. Pick \(\tilde{\imath}\) in that set and
\(\tilde{T}\subset I_1(\y)\setminus\{\tilde{\imath}\}\) containing
\(I_{2\lambda}^>(\x)\) with \(|\tilde{T}|=l-1\). Then
\(\tilde{T}\in\mathcal{S}_{l-1}(\omega(\x))\) and
\(\tilde{\imath}\in I^=_{[l]}(\x)\setminus\tilde{T}\). The residual
\(I_1(\y)\setminus(\tilde{T}\cup\{\tilde{\imath}\})\) lies in
\(I_{2\lambda}^=(\x)\setminus(\tilde{T}\cup\{\tilde{\imath}\})\) and has
cardinality \(|I_1(\y)|-l\leq s'-l\). Hence there exists
\(\tilde{U}\subset I_{2\lambda}^=(\x)\setminus(\tilde{T}\cup\{\tilde{\imath}\})\)
with \(I_1(\y)\setminus(\tilde{T}\cup\{\tilde{\imath}\})\subset\tilde{U}\) and
\(|\tilde{U}|=\min\bigl\{\bigl|I_{2\lambda}^=(\x)\setminus(\tilde{T}\cup\{\tilde{\imath}\})\bigr|,s'-l\bigr\}\).
Therefore
\(\tilde{T}\cup\{\tilde{\imath}\}\subset I_1(\y)\subset\tilde{T}\cup\{\tilde{\imath}\}\cup\tilde{U}\).
As already observed, \(\y_{G_i} \in P_{D_i}\left(\x_{G_i}\right)\)
\(\forall i \in I_1(\y)\), so (a) also holds for \(\y\).

Now suppose \(\omega(\x)_{[l]}<2\lambda\). By (b) and \eqref{T2},
\(I_1(\U)=T\cup\{i\}\in\mathcal{S}_l(\omega(\x))\) has cardinality \(l\),
hence \(I_1(\U)\cap I_{[l]}^<(\x)=\emptyset\), and
\begin{equation}
\label{ubest2}
\sum_{i \in I_1(\U)}\left(\omega(\x)_i-2 \lambda\right)=\sum_{i \in I_1(\U)\cap I_{2\lambda}^\geq(\x) }\left(\omega(\x)_i-2 \lambda\right)+\sum_{i \in I_1(\U)\cap I_{2\lambda}^<(\x)\cap I_{[l]}^\geq(\x) }\left(\omega(\x)_i-2 \lambda\right).
\end{equation}

By (b) again, since \(T\) contains indices of a set of the \(l-1\) largest
entries of \(\omega(\x)\), one has \(I_{2\lambda}^>(\x)\subset I_1(\U)\), and
then
\begin{equation}
    \label{greaterind}
    |I_1(\U) \cap I_{2\lambda}^>(\x)|=|I_{2\lambda}^>(\x)|\geq |I_1(\y) \cap I_{2\lambda}^>(\x)|.
\end{equation}
The same reasoning yields \(I^=_{2\lambda}(\x)\subset I_1(\U)\) and therefore
\begin{equation}
    \label{greaterind2}
    |I_1(\U) \cap I_{2\lambda}^{=}(\x)|=|I_{2\lambda}^{=}(\x)|\geq |I_1(\y) \cap I_{2\lambda}^{=}(\x)|.
\end{equation}
In particular, from \eqref{greaterind} and positivity of the summands,
\begin{equation}
\label{2lineq}
\sum_{i \in I_1(\U)\cap I_{2\lambda}^>(\x)}\left(\omega(\x)_i-2 \lambda\right) \geq \sum_{i \in I_1(\y)\cap I_{2\lambda}^>(\x)}\left(\omega(\x)_i-2 \lambda\right).
\end{equation}

Since by construction in this case \(|I_1(\U)|= l\) (\(U=\emptyset\) in
\eqref{T2}) and \(|I_1(\y)|\geq l\) (\(\y\in C_{l,s}\)), one has
\begin{align*}
    |I_1(\U) \cap I_{2\lambda}^>(\x)|+|I_1(\y) \cap I_{2\lambda}^\leq(\x)|
    &\geq|I_1(\y) \cap I_{2\lambda}^>(\x)|+|I_1(\y) \cap I_{2\lambda}^\leq(\x)| \eqjust{greaterind}\\
    &=|I_1(\y)|
    \\&\geq |I_1(\U)|
    \\&=|I_1(\U) \cap I_{2\lambda}^>(\x)|+|I_1(\U)\cap I_{2\lambda}^\leq(\x)\cap I_{[l]}^\geq(\x)| ,
\end{align*}
or equivalently,
\begin{equation}
    \label{inddiff}
    |I_1(\y) \cap I_{2\lambda}^\leq(\x)|\geq |I_1(\U)\cap I_{2\lambda}^\leq(\x)\cap I_{[l]}^\geq(\x)|.
\end{equation}

Since
\(|I_1(\y)\cap I_{2\lambda}^<(\x)|=|\bigl(I_1(\y)\cap I_{2\lambda}^\leq(\x)\bigr)|-|\bigl(I_1(\y)\cap I_{2\lambda}^=(\x)\bigr)|\)
(and likewise for \(\U\)), the inequality above on the first term and
\eqref{greaterind2} on the second yield in particular
\begin{equation}
    \label{inddiff2}
    |I_1(\y) \cap I_{2\lambda}^<(\x)|\geq |I_1(\U)\cap I_{2\lambda}^<(\x)\cap I_{[l]}^\geq(\x)|.
\end{equation}
That is, \(I_1(\y)\) may have to include indices from \(I_{[l]}^\leq(\x)\),
because
\(I_1(\U)\cap I_{2\lambda}^<(\x)\cap I_{[l]}^>(\x)=I_{2\lambda}^<(\x)\cap I_{[l]}^>(\x)\)
by construction (\(T\in\mathcal{S}_{l-1}(\omega (\x))\) and
\(T\subset I_1(\U)\)). One therefore obtains
\begin{equation}
\label{full}
\sum_{i \in I_1(\U)\cap I_{2\lambda}^<(\x)\cap I_{[l]}^\geq(\x)}\left(\omega(\x)_i-2 \lambda\right)
\geq \sum_{i \in I_1(\y) \cap I_{2\lambda}^<(\x)}\left(\omega(\x)_i-2 \lambda\right)
\eqjust{inddiff2},
\end{equation}
where it was also used that
\(\omega(\x)_{i}\geq \omega(\x)_j\ \forall i\in I^\geq_{[l]}(\x),\ j\in I^\leq_{[l]}(\x)\).

One may now form the chain of inequalities
\begin{equation}\label{lastchain}
\begin{aligned}
0 & \geq \sum_{i \in I_1(\y)\cap I_{2 \lambda}^>(\x)}\left(\omega(\x)_i-2 \lambda\right)-\sum_{i \in I_1(\U)\cap I_{2 \lambda}^>(\x)}\left(\omega(\x)_i-2 \lambda\right)
\eqjust{2lineq} \\
& =\sum_{i \in I_1(\y)\cap I_{2 \lambda}^\geq(\x)}\left(\omega(\x)_i-2 \lambda\right)-\sum_{i \in I_1(\U)\cap I_{2 \lambda}^\geq(\x)}\left(\omega(\x)_i-2 \lambda\right) \\
& \geq \sum_{i \in I_1(\y)\cap I_{2 \lambda}^\geq(\x)}\left(\omega(\x)_i-2 \lambda\right)+\sum_{i \in I_1(\y) \cap I^<_{2 \lambda}(\x)}\left(\omega(\x)_i-2 \lambda\right)
\eqjust{full} \\
& \qquad-\sum_{i \in I_1(\U)\cap I_{2 \lambda}^\geq(\x)}\left(\omega(\x)_i-2 \lambda\right)- \sum_{i \in I_1(\U)\cap I_{2\lambda}^<(\x)\cap I_{[l]}^\geq(\x)}\left(\omega(\x)_i-2 \lambda\right)\\
& =\sum_{i \in I_1(\y)\cap I_{2 \lambda}^\geq(\x)}\left(\omega(\x)_i-2 \lambda\right)+\sum_{i \in I_1(\y) \cap I^<_{2 \lambda}(\x)}\left(\omega(\x)_i-2 \lambda\right)
\eqjust{ubest2} \\
& \qquad-\sum_{i \in I_1(\U)}\left(\omega(\x)_i-2 \lambda\right) \\
& \geq 0. \eqjust{main1}
\end{aligned}
\end{equation}

Hence the chain holds with equality, which implies, for the second term in
particular,
\begin{equation}
    \label{greatereq}
\sum_{i \in I_1(\y)\cap I_{2 \lambda}^>(\x)}\left(\omega(\x)_i-2 \lambda\right)=\sum_{i \in I_1(\U)\cap I_{2 \lambda}^>(\x)}\left(\omega(\x)_i-2 \lambda\right),
\end{equation}
and consequently, since \(I_{2\lambda}^>(\x)\subset I_1(\U)\),
\begin{equation}\label{greatereq2}
    I_1(\y)\cap I_{2 \lambda}^>(\x)=I_1(\U)\cap I_{2 \lambda}^>(\x).
\end{equation}

For the last term of \eqref{lastchain}, \eqref{greatereq} also yields
(omitting the vanishing summands on \(I_{2\lambda}^=(\x)\))
\begin{equation}
\label{forcedeq}
\sum_{i \in I_1(\U)\cap I_{2\lambda}^<(\x)\cap I_{[l]}^\geq(\x)}\left(\omega(\x)_i-2 \lambda\right)= \sum_{i \in I_1(\y) \cap I_{2\lambda}^<(\x)}\left(\omega(\x)_i-2 \lambda\right).
\end{equation}
Since
\(I_1(\U)\cap I_{2\lambda}^<(\x)\cap I_{[l]}^>(\x)=I_{2\lambda}^<(\x)\cap I_{[l]}^>(\x)\),
the case
\(|I_1(\y) \cap I_{2\lambda}^<(\x)|> |I_1(\U)\cap I_{2\lambda}^<(\x)\cap I_{[l]}^\geq(\x)|\)
would imply
\(|I_1(\y) \cap I_{[l]}^\leq(\x)|>|I_1(\U) \cap I_{[l]}^=(\x)|\), that is,
\(\y\) would need smaller elements than \(\U\) in the sum in order to obtain
\(|I_1(\y) \cap I_{2\lambda}^<(\x)|> |I_1(\U)\cap I_{2\lambda}^<(\x)\cap I_{[l]}^\geq(\x)|\).
In that event \eqref{full} would hold as a strict inequality. By contradiction
with the equality \eqref{forcedeq}, one concludes
\(|I_1(\y) \cap I_{2\lambda}^<(\x)|= |I_1(\U)\cap I_{2\lambda}^<(\x)\cap I_{[l]}^\geq(\x)|\).
Again from \eqref{forcedeq}, this implies
$$
\sum_{i \in I_1(\y) \cap I^<_{[l]}(\x)}\left(\omega(\x)_i-2 \lambda\right)=0.
$$
Moreover, the latter relation and the negativity of the summands on
\(I_1(\y) \cap I^<_{[l]}(\x)\) guarantee
\begin{equation}\label{ylesser}
    I_1(\y) \cap I^<_{[l]}(\x)=\emptyset.
\end{equation}
In view of the last display, \eqref{greatereq2} and \(|I_1(\y)|\geq l\), and
recalling that \(I^=_{2\lambda}(\x)\subset I_1(\U)\), one is left with
\begin{equation}\label{yequals}
\begin{aligned}
|I_1(\y) \cap I^=_{2\lambda}(\x)|
&= |I_1(\U) \cap I^=_{2\lambda}(\x)|
= |I^=_{2\lambda}(\x)|, \\
I_1(\y) \cap I^=_{2\lambda}(\x)
&= I_1(\U) \cap I^=_{2\lambda}(\x)
= I^=_{2\lambda}(\x).
\end{aligned}
\end{equation}

Using \eqref{greatereq2}, \eqref{ylesser} and \eqref{yequals}, the sum over
\(I_1(\y)\) reduces to
$$
\sum_{i \in I_1(\y) }\left(\omega(\x)_i-2 \lambda\right)=\sum_{i \in I_1(\y)\cap I_{2\lambda}^\geq(\x) }\left(\omega(\x)_i-2 \lambda\right)+\sum_{i \in I_1(\y)\cap I_{2\lambda}^<(\x)\cap I_{[l]}^\geq(\x) }\left(\omega(\x)_i-2 \lambda\right),
$$
which finally produces the chain of equalities
\begin{align*}
\sum_{i \in I_1(\y) }\left(\omega(\x)_i-2 \lambda\right)&=
    \sum_{i \in I_1(\y)\cap I_{2\lambda}^\geq(\x) }\left(\omega(\x)_i-2 \lambda\right)+\sum_{i \in I_1(\y)\cap I_{2\lambda}^<(\x)\cap I_{[l]}^\geq(\x) }\left(\omega(\x)_i-2 \lambda\right)\\
    &=\sum_{i \in I_1(\U)\cap I_{2\lambda}^\geq(\x) }\left(\omega(\x)_i-2 \lambda\right)+\sum_{i \in I_1(\y)\cap I_{2\lambda}^<(\x)\cap I_{[l]}^\geq(\x) }\left(\omega(\x)_i-2 \lambda\right)
    \eqjust{greatereq}\\
    &=\sum_{i \in I_1(\U)\cap I_{2\lambda}^\geq(\x) }\left(\omega(\x)_i-2 \lambda\right)+\sum_{i \in I_1(\U)\cap I_{2\lambda}^<(\x)\cap I_{[l]}^\geq(\x) }\left(\omega(\x)_i-2 \lambda\right)
    \eqjust{forcedeq}\\
    &=\sum_{i \in I_1(\U) }\left(\omega(\x)_i-2 \lambda\right). \eqjust{ubest2}
\end{align*}

Finally, \eqref{proxy} also holds as an equality, which implies
\(\U \in \prox_h(\x)\).
It remains to exhibit \(\tilde{T}\in\mathcal{S}_{l-1}(\omega(\x))\),
\(\tilde{\imath}\in I^=_{[l]}(\x)\setminus\tilde{T}\) and
\(\tilde{U}=\emptyset\) such that \eqref{T2} holds with \(\U\) replaced by
\(\y\). From \eqref{greatereq2}, \eqref{ylesser}
and \eqref{yequals},
\(I_{2\lambda}^>(\x)\cup I_{2\lambda}^=(\x)\subset I_1(\y)\subset I_{[l]}^\geq(\x)\)
and \(|I_1(\y)|\geq l\). If \(|I_1(\y)|>l\), some index of \(I_1(\y)\) would have
\(\omega(\x)_i-2\lambda<0\) and could be removed, contradicting equality of
the sums with \(I_1(\U)\). Hence \(|I_1(\y)|=l\). If there existed
\(j\in I_{[l]}^>(\x)\setminus I_1(\y)\), some \(k\in I_1(\y)\cap I^=_{[l]}(\x)\) would
satisfy \(\omega(\x)_k=\omega(\x)_{[l]}<\omega(\x)_j\); the swap would
increase the sum, again a contradiction. Hence \(I_{[l]}^>(\x)\subset I_1(\y)\) and
\(I_1(\y)\in\mathcal{S}_l(\omega(\x))\). Since \(|I_{[l]}^>(\x)|\leq l-1\), the set
\(I_1(\y)\cap I^=_{[l]}(\x)\) is nonempty. Pick \(\tilde{\imath}\) in that set,
\(\tilde{T}\coloneqq I_1(\y)\setminus\{\tilde{\imath}\}\) and
\(\tilde{U}\coloneqq\emptyset\). Then
\(\tilde{T}\in\mathcal{S}_{l-1}(\omega(\x))\),
\(\tilde{\imath}\in I^=_{[l]}(\x)\setminus\tilde{T}\) and
\(I_1(\y)=\tilde{T}\cup\{\tilde{\imath}\}\).
As already discussed, \(\y\) also satisfies (a), which completes the proof.
\end{proof}

\section{Closed form and a computational scheme}
\label{sec:closed}

The previous result yields an explicit description of \(\prox_h\).
If \(|I_{\mathcal{A}}(\x)|<l\), then \(\prox_h(\x)=\emptyset\). Otherwise,
$$
\prox_h(\x)=\left\{\mathbf{U}_{\mathcal{A}(T)} \y\mid \y \in P_{\mathcal{B}(T)}\left(\x_{\mathcal{A}(T)}\right),\ T\in\mathcal{T}(\x)\right\},
$$
where
$$
\mathcal{T}(\x)=
\begin{cases}\left\{T\mid I_{[s']}^>({\x}) \subsetneq T \subset I_{[s']}^\geq({\x}),\ |T|=s'\right\}, &\ \omega(\x)_{[l]}>2 \lambda,\ \omega(\x)_{[s']}>2 \lambda \\
\left\{T\mid I_{[s']}^>({\x}) \subset T \subset I_{[s']}^\geq({\x}),\ |T| \leq s'\right\}, &\ \omega(\x)_{[l]}>2 \lambda,\ \omega(\x)_{[s']}=2 \lambda \\
\left\{T\mid I_{2\lambda}^>({\x}) \subset T \subset I_{2\lambda}^\geq({\x}),\ |T| < s'\right\}, &\ \omega(\x)_{[l]}>2 \lambda,\ \omega(\x)_{[s']}<2 \lambda
\\\left\{T\mid I_{[l]}^>({\x}) \subsetneq T \subset I_{[l]}^\geq({\x}),\ l\leq|T| \leq s'\right\}, &\ \omega(\x)_{[l]}=2 \lambda  \\
\left\{T\mid I_{[l]}^>({\x}) \subsetneq T \subset I_{[l]}^\geq({\x}),\ |T|=l\right\}, &\ \omega(\x)_{[l]}<2 \lambda
\end{cases}.
$$
This description suggests the following procedure for computing one element of
\(\prox_h\). The algorithm evaluates \(I_{\mathcal{A}}(\x)\), returns
\(\emptyset\) if fewer than \(l\) groups are activatable, and selects
\(\mathcal{S}_{s'}(\omega(\x))\) or \(\mathcal{S}_l(\omega(\x))\).

\begin{algorithm}[ht]
\caption{Group-sparse proximal mapping}
\label{alg:proxh}
\Input{\(\x\in\R^n\),\ \(l\leq s\leq m\) and \(\lambda \in \R\)}\\
\Output{\(\U\in\prox_h(\x)\) or \(\emptyset\)}
    \begin{algorithmic}[1]
    \State compute \(I_{\mathcal{A}}(\x)\)
    \If{\(l\geq 1\) \textbf{and} \(|I_{\mathcal{A}}(\x)|<l\)}
        \State \Return \(\emptyset\)
    \EndIf
    \State \(s'\gets\min\{s,|I_{\mathcal{A}}(\x)|\}\)
    \State compute \(S\in \mathcal{S}_{s'}(\omega(\x))\) (arbitrarily chosen)
\If{\(\omega(\x)_{[l]} >2\lambda\)}
\State \(T\gets\{i\in S\mid\omega(\x)_i >2\lambda\}\)
\Else
\State use \(S\) to compute \(T\in \mathcal{S}_l(\omega(\x))\)
\EndIf
\State \(\U\gets \mathbf{U}_{\mathcal{A}(T)} P_{\mathcal{B}(T)}\left(\x_{\mathcal{A}(T)}\right)\)
\end{algorithmic}
\end{algorithm}

\begin{remark}
If \(\omega(\x)_{[l]} \leq 2\lambda\), it is enough to compute a selection
\(T\in\mathcal{S}_l(\omega(\x))\) of \(l\) of the largest entries of
\(\omega(\x)\), which is cheaper than an initial selection
\(S\in\mathcal{S}_s(\omega(\x))\) when \(l<s\). Moreover, entries with
\(\omega(\x)_i \leq 2\lambda\) are irrelevant when
\(\omega(\x)_{[l]} > 2\lambda\), so \(S\) can sometimes be shortened. A more
refined implementation would first obtain a selection of \(l\) elements, read
off \(\omega(\x)_{[l]}\), and stop if \(\omega(\x)_{[l]} \leq 2\lambda\);
otherwise one continues until \(S\in \mathcal{S}_s(\omega(\x))\) is filled or
a newly inserted entry is at most \(2\lambda\).

When evaluating \(\omega(\x)\) in \eqref{omega}, the projection onto each
\(D_i\) computed for \(d_{D_i}\) should be stored and reused in
\(P_{\mathcal{B}(T)}\).
\end{remark}

\subsection{Circumventing an empty \(\prox\)}
\label{sec:emptyfix}

\Cref{alg:proxh} correctly returns \(\emptyset\) when \(l\ge 1\) and
\(|I_{\mathcal{A}}(\x)|<l\). That answer is faithful to \eqref{h}, but a
proximal-gradient loop cannot take a step. A direct fix is to cut from each \(D_i\) the open ball of radius
\(\varepsilon_i>0\) about the origin, in the Euclidean or \(\ell_\infty\)
norm,
\[
D_i^{\varepsilon}\coloneqq D_i\setminus\bigl\{\z\in\R^{n_i}:\|\z\|<\varepsilon_i\bigr\}.
\]
The origin remains available
through \(B\), but an active group must leave that ball. If \(\varepsilon_i\)
is small enough that \(D_i\) is not contained in the ball, then
\(D_i^{\varepsilon}\) is closed and does not contain \(\mathbf{0}\), so every
remaining group is activatable for every \(\x\). The modified regularizer
\(h^{\varepsilon}\) then has a nonempty \(\prox_{h^{\varepsilon}}\) everywhere. The radius can come from the application, for
instance a numerical floor such as \(10^{-8}\) times a typical group scale.

If emptiness is only the coincidence
\(P_{D_i}(\x_{G_i})=\{\mathbf{0}\}\) at a single point, a tiny perturbation of
\(\x_{G_i}\) is enough as a numerical safeguard.

\section{Concluding remarks}
\label{sec:end}

We have extended the proximal calculus of \cite{beck2018} from an upper
cardinality bound to a genuine window \(l\le\|g(\x)\|_0\le s\). The lower
bound changes the geometry of \(\prox_h\). It can be empty, and when
\(\omega(\x)_{[l]}\le 2\lambda\) one must activate zero-gain (or even
negative-gain) groups in order to meet the quota \(l\). Both phenomena are
invisible when \(l=0\).

The resulting oracle can be used in proximal-gradient schemes for
\eqref{P}. If \(|I_{\mathcal{A}}(\x)|<l\), the hole in \(D_i\) of
\cref{sec:emptyfix} restores a nonempty \(\prox\) to which the theorem applies.
We do not recast the stationarity hierarchy of
\cite{beck2018} for the window \(C_{l,s}\). Once \(\prox_h\) is available those
conditions transfer formally, except that proximal stationarity is undefined
at centres with fewer than \(l\) activatable groups. Natural next steps are a
numerical study of the role of \(l\) on applications such as
\cref{ex:carteira,ex:imbuia}, and the use of this \(\prox\)
inside forward--backward methods.

\bibliographystyle{plainnat}
\bibliography{paper_extra}

\end{document}
